\documentclass[10pt]{beamer}
\usetheme{Madrid}
\usecolortheme{seahorse}

\usepackage[T1]{fontenc}
\usepackage{lmodern}
\usepackage{listings}
\usepackage{xcolor}
\usepackage{amsmath}
\usepackage{booktabs}

\lstset{
  basicstyle=\ttfamily\scriptsize,
  keywordstyle=\color{blue},
  commentstyle=\color{gray},
  stringstyle=\color{purple},
  showstringspaces=false,
  breaklines=true,
  language=Python,
}

\setbeamertemplate{navigation symbols}{}
\setbeamertemplate{footline}[frame number]

\title[GWPCA in spml]{Geographically Weighted PCA for \texttt{spml}}
\subtitle{A suggested implementation}
\author{}
\institute{PySAL --- SDSU spatial-ML group}
\date{}

\begin{document}

\begin{frame}
  \titlepage
  \begin{center}
    \small A 10-minute walk-through of the design, the maths, and what is
    already on the \texttt{gwpca} branch.
  \end{center}
\end{frame}

% ---------------------------------------------------------
\begin{frame}{The request}
  \textbf{Extend the GW framework to unsupervised models,
  starting with GWPCA.}
  \begin{itemize}
    \item Today \texttt{spml} fits one \emph{supervised} scikit-learn
      estimator per location: \texttt{GWLinearRegression},
      \texttt{GWLogisticRegression}, GW random forests, GW gradient boosting.
    \item All operate on pairs $(X, y)$.
    \item We want a branch that works on the \textbf{feature matrix $X$
      alone} and decomposes its \emph{local} covariance structure.
    \item First estimator: \textbf{Geographically Weighted PCA}
      \footnotesize (Harris, Brunsdon \& Charlton, 2011).
  \end{itemize}
\end{frame}

% ---------------------------------------------------------
\begin{frame}{Recap: the ``map of models'' idea}
  For each location $i$:
  \begin{enumerate}
    \item Build spatial weights $w_{ij}$ from a \textbf{kernel} and a
      \textbf{bandwidth} (fixed distance or adaptive $k$-NN), via
      \texttt{libpysal.graph}.
    \item Fit a local model at $i$, weighting each observation $j$ by
      $w_{ij}$.
    \item Store the local parameters for $i$.
  \end{enumerate}
  \vspace{0.6em}
  GWPCA keeps steps 1 and 3 unchanged. Only \textbf{step 2} is different:
  no target, no \texttt{.fit(X, y)} --- a direct eigendecomposition instead.
\end{frame}

% ---------------------------------------------------------
\begin{frame}{What GWPCA produces}
  A principal components analysis whose loadings and variances
  \textbf{vary across space}:
  \begin{itemize}
    \item local loadings --- how variables combine into components
      \emph{here}
    \item local explained variance --- how much structure each component
      captures \emph{here}
    \item the \textbf{winning variable} --- which variable dominates
      component 1 at each location (the classic GWPCA map)
    \item proportion of total variation (PTV) retained by the first
      $k$ components
  \end{itemize}
  \vfill
  \centering
  \textit{One global rotation $\Rightarrow$ a surface of rotations.}
\end{frame}

% ---------------------------------------------------------
\begin{frame}{The one real wrinkle}
  \begin{block}{\texttt{sklearn.decomposition.PCA} has no \texttt{sample\_weight}}
    We cannot just hand weighted rows to scikit-learn.
  \end{block}
  \vspace{0.5em}
  So the local model is \textbf{not} an sklearn estimator fit --- it is a
  weighted covariance matrix built directly and eigendecomposed with NumPy.
  \begin{itemize}
    \item cheap: one $p \times p$ symmetric eigen-problem per location
    \item deterministic --- no solver, no \texttt{random\_state}
    \item the global baseline \emph{is} a plain \texttt{PCA}, for comparison
  \end{itemize}
\end{frame}

% ---------------------------------------------------------
\begin{frame}{The maths (per location $i$)}
  Geographically weighted mean of variable $k$:
  \[
    \bar{x}_k(i) \;=\; \frac{\sum_j w_{ij}\, x_{jk}}{\sum_j w_{ij}}
  \]
  Geographically weighted covariance:
  \[
    \Sigma_{k\ell}(i) \;=\;
      \frac{\sum_j w_{ij}\,\bigl(x_{jk}-\bar{x}_k(i)\bigr)
                          \bigl(x_{j\ell}-\bar{x}_\ell(i)\bigr)}
           {\sum_j w_{ij}}
  \]
  With \texttt{standardize=True} (the default), $x$ is first z-scored
  \textbf{globally}, once, with the population SD:
  \[
    x_{jk} \;\leftarrow\; \frac{x_{jk} - \bar{x}_k}{\operatorname{sd}(x_k)}
  \]
  \footnotesize This is \emph{not} a local correlation matrix: the diagonal of
  $\Sigma(i)$ is not forced to 1. It only becomes a correlation matrix in the
  large-bandwidth limit.
\end{frame}

% ---------------------------------------------------------
\begin{frame}{The maths --- decomposition and outputs}
  Symmetric eigendecomposition:
  \[
    \Sigma(i) \;=\; V(i)\,\Lambda(i)\,V(i)^{\!\top},
    \qquad \lambda_1(i) \ge \dots \ge \lambda_p(i)
  \]
  \begin{itemize}
    \item columns of $V(i)$ --- local loadings; $\Lambda(i)$ --- local
      variances
    \item explained-variance ratio: $\lambda_m(i) / \sum_m \lambda_m(i)$
    \item score for each neighbour $j$ of $i$ (raw, \emph{not}
      mean-centered --- matches \texttt{GWmodel}, more shortly):
      $s_j(i) = x_j^{\!\top} V(i)_{[:,\,1:k]}$
    \item winning variable: $\displaystyle \arg\max_k \bigl|V_{k1}(i)\bigr|$
  \end{itemize}
  \vspace{0.4em}
  \footnotesize Sign convention: force the largest-magnitude loading of
  each component positive (per location).
\end{frame}

% ---------------------------------------------------------
\begin{frame}{How it fits the existing architecture}
  \begin{center}
  \small
  \begin{tabular}{@{}ll@{}}
    \toprule
    \textbf{Existing piece} & \textbf{In GWPCA} \\
    \midrule
    \texttt{\_build\_weights} & reuse as-is (kernel graph, adaptive/fixed) \\
    parallel neighbourhood fan-out & reuse (drop the \texttt{\_y} column) \\
    \texttt{\_fit\_local} & replaced by \texttt{\_fit\_neighborhood} \\
    prediction helpers & reuse later for \texttt{transform} \\
    \texttt{keep\_models} (mem / disk) & reuse: store local mean + loadings \\
    IC / hat-matrix / AIC code & \textbf{skip} (\texttt{\_model\_type = "gwpca"}) \\
    \texttt{fit\_global\_model} & reuse: baseline is a plain \texttt{PCA} \\
    \texttt{BandwidthSearch} & extend later (CV reconstruction error) \\
    \bottomrule
  \end{tabular}
  \end{center}
  \vfill
  \centering
  \normalsize Net new code is small and localised in one module.
\end{frame}

% ---------------------------------------------------------
\begin{frame}[fragile]{New module: \texttt{spml/decomposition.py}}
\begin{lstlisting}
from spml.decomposition import GWPCA

gwpca = GWPCA(
    n_components=2,
    bandwidth=30,
    fixed=False,          # adaptive k-NN
    kernel="bisquare",
    standardize=True,     # global z-score + local cov (see R note)
    keep_models=True,     # needed for transform() later
).fit(X, geometry=gdf.representative_point())  # points only
\end{lstlisting}
  \begin{itemize}
    \item \texttt{fit(X, geometry)} --- no \texttt{y}; geometry must be
      points
    \item Same split as \texttt{BaseRegressor} $\to$
      \texttt{GWLinearRegression}: generic machinery in the base class,
      method-specific code in the subclass (next slide)
  \end{itemize}
\end{frame}

% ---------------------------------------------------------
\begin{frame}{Base class vs \texttt{GWPCA}}
  \begin{columns}[T]
    \begin{column}{0.5\textwidth}
      \textbf{\texttt{BaseGWDecomposition}}
      \begin{itemize}\small
        \item validation, global standardization
        \item weights, parallel loop, model storage
        \item \texttt{\_fit\_neighborhood}: weights + local mean, then calls
          \texttt{\_decompose\_local}
        \item \texttt{\_scale\_eigenvalues}: divide by each location's
          \emph{own} weight sum
        \item generic attributes: \texttt{components\_},
          \texttt{explained\_variance[\_ratio]\_}
        \item \texttt{model=} builds the global baseline
      \end{itemize}
    \end{column}
    \begin{column}{0.5\textwidth}
      \textbf{\texttt{GWPCA}}
      \begin{itemize}\small
        \item \texttt{\_\_init\_\_} with \texttt{model=PCA}
        \item \texttt{\_decompose\_local}: scatter matrix, \texttt{eigh},
          sign rule, neighbour scores
        \item \texttt{\_scale\_eigenvalues} override: the GWmodel bug
          (shared divisor)
        \item \texttt{fit} wrapper: \texttt{local\_scores\_},
          \texttt{local\_loadings\_}, PTV, winning variable
      \end{itemize}
    \end{column}
  \end{columns}
  \vfill
  \centering
  \footnotesize The GWmodel bug lives in one method: dropping it means
  deleting the override.
\end{frame}

% ---------------------------------------------------------
\begin{frame}[fragile]{Fitted attributes (scikit-learn flavoured)}
\begin{lstlisting}
gwpca.components_               # (n_samples, n_components, n_features)
gwpca.explained_variance_       # (n_samples, n_components)
gwpca.explained_variance_ratio_ # (n_samples, n_components)
gwpca.local_eigenvalues_        # all eigenvalues, descending
gwpca.proportion_total_var_     # PTV of the first k components
gwpca.winning_variable_         # argmax |loading| on PC1
gwpca.local_scores_             # per-neighbour scores (see R note)
gwpca.local_loadings_           # tidy MultiIndex (location, component)
gwpca.global_model_             # sklearn PCA baseline
\end{lstlisting}
\end{frame}

% ---------------------------------------------------------
\begin{frame}{What is on the \texttt{gwpca} branch now}
  \textbf{Implemented (the unambiguous core):}
  \begin{itemize}
    \item weighted mean / covariance, global standardization, \texttt{eigh},
      descending sort, per-location sign convention
    \item all the attributes on the previous slide
    \item \texttt{graph=} and \texttt{geometry+bandwidth} paths; fixed and
      adaptive; \texttt{keep\_models}; global baseline
    \item \textbf{19 tests}; rebased on current \texttt{main}, full suite
      green (405 passed, 3 skipped)
    \item user-guide notebook (Guerry data), \texttt{api.rst}, references
  \end{itemize}
  \textbf{Verified against ground truth:}
  \begin{itemize}
    \item matches a hand-rolled weighted covariance $+$ \texttt{eigh}
    \item \texttt{boxcar} kernel $+$ huge bandwidth reproduces global
      \texttt{sklearn} PCA loadings and ratios to machine precision
  \end{itemize}
  \vspace{0.3em}
  \footnotesize Since: code-reviewed and hardened --- bandwidth input
  validation (matching the supervised GW models), dead-code/lint cleanup,
  vectorized attribute assembly, doc fixes. All six review findings closed.
\end{frame}

% ---------------------------------------------------------
\begin{frame}{A wrinkle found while cross-checking against R}
  Went a step further than the checks on the previous slide: ran real
  \texttt{GWmodel::gwpca} (R) on identical data and diffed the output
  directly --- and found a genuine bug in the reference implementation.
  \begin{block}{The bug}
    \texttt{gwpca()} normalises every location's reported eigenvalues by
    \texttt{sum(wt)} taken from whichever neighbourhood was \emph{last
    processed} in the input's row order (a stray/reused loop variable in the
    R source) --- \textbf{not} that location's own neighbourhood weight-sum.
  \end{block}
  \begin{itemize}
    \item Harmless only when every neighbourhood has the same weight-sum,
      which in practice means an adaptive \texttt{boxcar} kernel.
    \item Wrong for any distance-decay kernel, \textbf{including the default
      adaptive bisquare} (Guerry, 30 neighbours: weight-sums range from 8.6 to
      13.9), and for fixed kernels over irregular points: eigenvalue
      \emph{magnitudes} are mis-scaled for every location except the one
      processed last.
    \item Loadings and \texttt{explained\_variance\_ratio\_}/PTV are
      \textbf{unaffected} --- the shared divisor cancels out of any
      within-row ratio and never touches the eigenvectors.
  \end{itemize}
\end{frame}

% ---------------------------------------------------------
\begin{frame}[fragile]{Our call: replicate it, on purpose}
  \texttt{spml} now deliberately reproduces this exact behaviour, for
  numerical parity with the reference implementation:
  \begin{itemize}
    \item \texttt{local\_eigenvalues\_} / \texttt{explained\_variance\_}
      share one divisor (the last-processed neighbourhood's weight-sum)
      $\Rightarrow$ now \textbf{depend on \texttt{X}'s row order} at fit time.
    \item Documented in \texttt{decomposition.py}'s module docstring
      (``Deliberate parity with \texttt{GWmodel::gwpca}'') and locked in by
      \texttt{test\_gwpca\_eigenvalue\_scale\_depends\_on\_row\_order}.
    \item \texttt{test\_gwpca\_matches\_r\_gwmodel\_reference}: eigenvalues,
      loadings and per-neighbour scores match real \texttt{GWmodel::gwpca}
      output.
    \item \texttt{gwpca()} has no \texttt{scale.vars} argument; users call
      \texttt{scale()} first. \texttt{standardize=True} does the same with
      the population SD: same loadings, eigenvalues differ by $(n-1)/n$.
    \item ``Last processed'' is looked up by \texttt{X}'s last row label, not
      the spatial graph's internal order.
  \end{itemize}
  \vfill
  \centering
  \footnotesize Worth reporting upstream to the \texttt{GWmodel} maintainer.
\end{frame}

% ---------------------------------------------------------
\begin{frame}[fragile]{Two more GWmodel bugs: bandwidth cross-validation}
  \texttt{gwpca.cv} (bandwidth score) and \texttt{gwpca.cv.contrib}
  (per-observation residuals, used for outliers) compute, for each left-out
  $i$:
\begin{lstlisting}[language=R]
v <- v %*% t(v)                               # projector V_k V_k^T
score <- score + sum((x[i, ] - x[i, ] %*% v))^2
\end{lstlisting}
  \begin{itemize}
    \item \textbf{Uncentred:} reconstructs raw $x_i$, while the loadings come
      from data centred on the local mean.
    \item \textbf{\texttt{sum(r)\^{}2}, not \texttt{sum(r\^{}2)}:} squares the
      \emph{sum} of the residuals, so positive and negative residuals cancel.
    \item On Guerry: the textbook score picks \textbf{14} neighbours, the
      GWmodel score \textbf{44}.
    \item Brunsdon's tutorial supplies a ``fixed'' \texttt{gwpca.cv.contrib},
      but it has the same bug.
  \end{itemize}
  \vfill
  \centering
  \footnotesize Unlike the eigenvalue bug, this changes \emph{which bandwidth
  and which outliers} are chosen, so the proposal is a correct score by
  default, with GWmodel's version as an option.
\end{frame}

% ---------------------------------------------------------
\begin{frame}{Open decisions --- need the API owners}
  \begin{itemize}
    \item \textbf{Base-class name} (\texttt{BaseGWDecomposition} or
      \texttt{BaseDecomposition}) and whether \texttt{model=} stays (it now
      builds the global baseline).
    \item \textbf{GWmodel eigenvalue bug:} keep it as the default, put it
      behind an option, or drop it.
    \item \textbf{\texttt{local\_scores\_}:} keep GWmodel's per-neighbour,
      uncentred scores, or also offer one centred score per location.
    \item \textbf{CV score} for bandwidth selection: textbook or GWmodel.
    \item \textbf{Cross-location sign \& order} --- eigenvectors are defined
      up to sign, and component order can swap between neighbours.
      Per-location convention is in; \emph{aligning across space} (Procrustes
      rotation to the global solution) is still open --- and the usual
      critique of GWPCA maps.
    \item \textbf{Fixed vs adaptive \texttt{n\_components}} --- fixed for now.
  \end{itemize}
\end{frame}

% ---------------------------------------------------------
\begin{frame}{Roadmap}
  \textbf{First PR (this branch):} \texttt{BaseGWDecomposition} $+$
  \texttt{GWPCA} $+$ core attributes $+$ tests $+$ user-guide notebook,
  \texttt{api.rst} and references. \emph{Code done; PR to
  \texttt{pysal/spml} not yet opened.}
  \vspace{0.6em}

  \textbf{Follow-up PRs:}
  \begin{enumerate}
    \item \texttt{BandwidthSearch} integration --- leave-one-out CV
      (the notebook computes it by hand for now)
    \item Monte Carlo test of local variation
    \item cross-location alignment (Procrustes)
    \item \texttt{transform()} on new data --- nearest / ensemble local
      models (reuse the prediction helpers)
    \item \texttt{batch\_size} for large samples; robust GWPCA and outliers
  \end{enumerate}
\end{frame}

% ---------------------------------------------------------
\begin{frame}{Get involved}
  \begin{itemize}
    \item Branch: \texttt{gwpca}; upstream PR to \texttt{pysal/spml} in
      preparation
    \item Start with the user-guide notebook,
      \texttt{docs/source/gwpca.ipynb}
    \item Good pickups for the group: bandwidth search (CV), the Monte Carlo
      test, the alignment question, \texttt{transform}
  \end{itemize}
  \vfill
  \centering
  \Large \texttt{from spml.decomposition import GWPCA}
\end{frame}

\end{document}
